\section{Discussion}
\label{sec:discussion}

\subsection{Expected versus actual results}
\label{sec:discussion-expected}

We pre-registered three specific expectations and observed three outcomes that differed from them, in each case along a pre-committed path.

\textbf{Convergence at GATE-2.} The Stage 2 protocol at 25 rounds had observed $0/15$ strict-convergent repetitions at $n = 2$. We extended the horizon to 50 rounds in Stage 2b on the hypothesis that convergence at $n = 2$ was horizon-limited. The observed result is $1/15$ ($6.7\%$): a marginal improvement that does not lift convergence at $n = 2$ above the pre-registered Stage 3 PROCEED threshold. The duopoly is not horizon-limited under the current protocol; smaller-agent-count duopolies appear to occupy a wider basin within which terminal-state characterization requires longer horizons \emph{and} higher repetition counts. Stage 3 expansion to 100-round horizons at $n_{\text{rep}} \geq 30$ is the pre-registered fix.

\textbf{Basin count at GATE-5.} Prior exploratory work at the 2026-04-21 panel synthesis examined a $k = 4$ quadrimodal cluster structure at fixed $k$. The formal Addendum \#1 \S A bootstrap procedure recovers $k = 2$ in $7/10$ resamples (cleared by exactly one resample; supplementary 100-resample $\times$ 10-seed sweep returns $72.0\%$ aggregated and $10/10$ modal-winner consistency). The pre-committed downgrade clause (branch (c)) fires literally on the data: \emph{retract the quadrimodal narrative entirely and report bimodal throughout, which strengthens alignment with} \citep{Cao_2026_llm_collusion}. We honor the pre-registered branch and report $k = 2$ throughout this paper; the exploratory $k = 4$ centroids are retained in the supplementary materials but do not support any preprint claim.

\textbf{Reasoning-length as within-basin signature.} The 2026-04-21 panel synthesis had referenced an exploratory $r(\text{reasoning}, \Delta_{\text{profit}}) = -0.78$ at GATE-5 under an all-rounds aggregation. The locked Addendum \#1 \S C window (per-rep means over rounds 41--50) returns $r = -0.30$ with Fisher-$z$ 95\% CI $[-0.60, +0.06]$, which spans zero. The basin fixed-effects regression returns $\hat{\beta}_{\text{reasoning}} = +6.25 \times 10^{-4}$ with $p = 0.665$ at GATE-5; the locked decision rule fires that reasoning-length is not a within-basin signature. Under the strict effect-size + confidence-interval test on the cluster-mean difference at the formally selected $k = 2$, the between-basin pattern is descriptively present with Cohen's $d = 0.60$ medium effect, but Welch's $p = 0.110$ and bootstrap BCa CI $[-6.6, +66.7]$ chars span zero --- we therefore retract the stronger ``between-basin discriminator'' framing and report the descriptively-present medium-effect pattern with deferral to Stage 4 cross-model replication for formal confirmation.

In each case the post-formalization conclusion is more conservative than the pre-formalization framing, and in each case the downgrade is data-driven through a pre-committed branch, not retroactively rationalized.

\subsection{Why the asymmetry: basin width and agent count}
\label{sec:discussion-asymmetry}

The headline asymmetry --- comparable $\Delta_{\text{profit}}$ at convergence ($0.49$ at $n = 5$, $0.45$ at $n = 2$) but sharply different convergence rates ($30/30$ versus $1/15$) --- is consistent with a wider basin of attraction at smaller agent counts. We sketch the mechanistic argument and the falsifiable prediction it generates.

The basin-of-attraction framework \citep{wunder_littman_babes_2010_qlearning_basin,benaim_hofbauer_sorin_2005_stochastic_approximation_differential_inclusions,askenazi2024bounds} treats terminal-state distributions as samples from the joint structure of (a) noise in agent learning dynamics --- modulated by initialization and exploration rate --- and (b) the topology of the value-function landscape. At higher effective noise, trajectories take longer to commit to a basin, exhibit larger within-basin drift, and a larger fraction of trajectories occupy non-stable transitional states at any fixed horizon. We adopt this framework operationally; agent count modulates effective in-context exploration for LLM agents analogously to how exploration rate modulates basin width for $Q$-learners \citep{askenazi2024bounds}. Whether ``basin width'' transports literally from parametric $Q$-learner theory to LLM in-context inference is a Stage 4 question.

For LLM agents, ``noise-to-signal ratio'' is not parameterized by an exploration schedule but is implicit in the in-context reasoning process: the same prompt with the same prior history can produce a different price commitment, mediated through stochastic decoding and through the variable-depth chain-of-thought trace the agent constructs. Reducing the agent count from $n = 5$ to $n = 2$ removes structural features that anchor the cross-agent equilibrium-selection signal: with five agents, the price vector itself contains substantial cross-sectional information that constrains each agent's update; with two agents, the cross-sectional signal collapses and per-agent decoding noise dominates. Under this reading, the duopoly $n = 2$ exhibits a wider basin in which more trajectories occupy long transients that have not equilibrated by round 50.

The data are consistent with this reading at multiple levels. At GATE-5, $30/30$ trajectories satisfy the strict 1\%-over-5-rounds convergence criterion within 50 rounds, $30/30$ are basin-stable across the rounds-31--40 vs rounds-41--50 windows, and the bimodal $k = 2$ structure is recovered in $7/10$ resamples (and $72.0\%$ aggregated across the supplementary 1000-resample sweep) by the gap statistic --- a tight, well-discriminated terminal-state distribution. At GATE-2, only $1/15$ trajectories satisfy the strict criterion, $12/15$ are basin-stable but with Wilson 95\% CI overlapping the $0.50$ baseline, and the $k = 1$ unimodal recovery in $10/10$ resamples is itself a finding: the duopoly $\Delta_{\text{profit}}$ distribution is not multimodal at $n_{\text{rep}} = 15$, which is consistent with a wider single basin rather than with multiple narrow basins.

The closest cousin signal is \citet{BrownMackay_2024_nber_pricing_cadence}'s cadence-asymmetry result: asymmetric dynamic structure between two pricing algorithms produces the same collusive outcome Calvano predicts in symmetric simulation. Our finding extends the asymmetry frame from cross-firm cadence to cross-condition agent-count: the same $\Delta_{\text{profit}}$ landing point can be reached through structurally different convergence paths.

\paragraph{Falsifiable prediction.}
The basin-width-by-agent-count hypothesis generates a concrete prediction: at Stage 3 with $n_{\text{rep}} \geq 30$ and horizon $\geq 100$ rounds, the strict-convergence rate at $n = 2$ should remain materially below the corresponding rate at $n = 5$ (i.e., the basin-width gap is structural, not a pure sample-size or horizon artifact). If Stage 3 instead recovers comparable convergence rates at the two agent counts, the basin-width hypothesis is rejected and the asymmetry must be attributed to horizon or sample-size effects alone. We pre-commit this prediction here.

\paragraph{Tension with canonical IO theory.}
This basin-width-by-agent-count framing reverses a canonical industrial-organization prediction: classical oligopoly theory \citep{stigler_1964_theory_of_oligopoly} attributes collusion difficulty to the free-rider problem, which grows monotonically with the number of competing firms --- implying \emph{narrower} collusive basins at higher $n$, not wider. We acknowledge this reversal directly. Two readings are consistent with the Stage 2b data, and only the Stage 3 expansion plus the cross-family Stage 4 design discriminate between them: \emph{(a)} the canonical free-rider mechanism applies to LLM agents but is dominated by an in-context cross-sectional signal effect --- with $n=5$ agents the per-round price vector contains substantially more information than at $n=2$, and the LLM's natural-language reading of that vector functions as an equilibrium-selection cue absent from the $Q$-table updating in classical algorithmic-pricing models; or \emph{(b)} the apparent convergence-rate asymmetry is an artifact of horizon (50 rounds insufficient to characterize the duopoly basin) and sample size ($n_{\text{rep}} = 15$ at GATE-2 underpowered) that Stage 3 will dissolve. Reading \emph{(a)} is the basin-width-by-agent-count hypothesis; reading \emph{(b)} is the IO canon. We pre-committed reading \emph{(a)} as the hypothesis under test and reading \emph{(b)} as the null. The pre-committed Addendum \#5 anchor empirically reversed reading \emph{(a)} as initially formulated (\S\ref{sec:discussion-asymmetry-cot}) and the pre-committed Addendum \#6 talk-volume-controlled residualization established that the reversal survives the most plausible alternative explanation, supporting a refined reading \emph{(a$'$)} on Stage 2b data; the Stage 3 falsification clauses (\S\ref{sec:discussion-falsifiability}) jointly resolve reading \emph{(a$'$)} versus reading \emph{(b)}. We do not claim that refined reading \emph{(a$'$)} dominates Stigler's free-rider mechanism: \citet{stigler_1964_theory_of_oligopoly}'s prediction lives in the price-theoretic / payoff-mechanism regime, while reading \emph{(a$'$)} lives in the learning-dynamics + cross-sectional-information regime. The two operate in different ontologies; we claim \emph{(a$'$)} is the operative mechanism in the regime our protocol probes, not that it adjudicates against Stigler's mechanism in the regime Stigler describes.

\paragraph{Empirical test of reading (a) — Addendum \#5 (exploratory, RFC 3161 stamped 2026-05-03).}
We pre-committed Addendum \#5 to add a text-based empirical test of reading \emph{(a)}: a regex-based count of cross-agent references in chain-of-thought, normalized per other-agent (\S\ref{sec:discussion-asymmetry-cot}). The result is a directional reversal of reading \emph{(a)} as initially formulated: the per-other-agent CoT cross-reference rate is $0.508$ per round at $n = 5$ (95\% bootstrap CI $[0.500, 0.516]$) and $2.954$ per round at $n = 2$ (95\% bootstrap CI $[2.865, 3.045]$), CIs disjoint by $\sim 5\times$ with GATE-2 \emph{higher} --- the LLM at $n = 2$ allocates substantially more per-peer attention than at $n = 5$. The most plausible alternative explanation (talk-volume artifact: GATE-2 produces more total CoT text per round, scaling the reference count mechanically) was tested via the pre-committed Addendum \#6 talk-volume-controlled residualization (RFC 3161 stamped 2026-05-04, \S\ref{sec:discussion-asymmetry-cot}) and ruled out: mean total word counts are nearly identical across conditions ($\sim 195$ words/CoT-round in both), and the per-other-agent rate gap survives pooled OLS residualization with disjoint $95\%$ CIs. Refined reading \emph{(a$'$)} --- per-peer attention density at $n = 2$ is denser even controlling for talk volume --- is empirically supported on Stage 2b data. The Stage 4 cross-model question (Sonnet/Opus generalization vs Haiku-class artifact) remains; the talk-volume-confound question is closed.

\subsubsection{Cross-agent CoT references — Addendum \#5 result (exploratory)}
\label{sec:discussion-asymmetry-cot}

Per Addendum \#5 (RFC 3161 stamped 2026-05-03; SHA-256 \texttt{033cabda5fee0ee8c74d178ea2d69388e1dd5dfdf5084636c8860949c03eeedf}), we computed the per-CoT-round cross-agent reference rate at GATE-5 vs GATE-2 as a text-based empirical anchor for the §5.2 reading \emph{(a)} mechanism. The analyzer (\texttt{pilot/analyze-cot-cross-references.ts}, Bun runtime) counts (i) literal named-agent references (``Agent N'', ``competitor'', ``other agent/player/firm/seller'', ``peer'') within a $\pm 5$-token window of any price-relevant token (price, profit, margin, undercut, cooperation, monopoly, Nash, equilibrium, or numeric \$1.x), and (ii) co-referent references (``their'', ``the other'', ``the higher/lower'', ``across the market'', ``the market'') under the same window. Self-references are excluded by sentence-prefix filter (``I'', ``My'', ``We'', ``Our'') and $\pm 3$-token first-person window filter. The primary metric is the per-other-agent normalized rate (raw count divided by $N - 1$, where $N$ is the agent count), which isolates per-peer signal density from the trivial agent-count effect. Bootstrap 95\% CIs are computed at $B = 10{,}000$ resamples with fixed seed = 42.

\textbf{Result.} The decision rule fires the REVERSAL branch (Addendum \#5 \S B.3): the primary per-other-agent rate is disjoint between conditions with GATE-2 higher. Per-condition results in Table~\ref{tab:cot-cross-references}.

\begin{table}[h]
\centering
\caption{Cross-agent CoT reference rates per Addendum \#5. Both raw channels (literal, co-referent) and the primary per-other-agent normalized rate favor GATE-2; the primary rate at GATE-2 is approximately $5.8\times$ the rate at GATE-5, with disjoint 95\% bootstrap CIs.}
\label{tab:cot-cross-references}
\begin{tabular}{lccc}
\toprule
Channel & GATE-5 ($n_{\text{round}} = 7500$) & GATE-2 ($n_{\text{round}} = 1500$) & CIs \\
\midrule
Literal raw refs / round & $1.579$ $[1.550, 1.607]$ & $1.753$ $[1.687, 1.819]$ & disjoint \\
Co-referent raw refs / round & $0.454$ $[0.440, 0.468]$ & $1.201$ $[1.134, 1.269]$ & disjoint \\
\textbf{Per-other-agent rate} & $\mathbf{0.508}$ $\mathbf{[0.500, 0.516]}$ & $\mathbf{2.954}$ $\mathbf{[2.865, 3.045]}$ & \textbf{disjoint} \\
\bottomrule
\end{tabular}
\end{table}

The per-other-agent rate at $n = 2$ is approximately $5.8\times$ the rate at $n = 5$ under per-other-agent normalization (raw literal-channel ratio $1.11\times$; raw co-referent ratio $2.65\times$); $n_{\text{round}} = 7500$ at GATE-5 and $1500$ at GATE-2 reflect the $30 \times 5 \times 50$ and $15 \times 2 \times 50$ agent-round structure. Both conditions' raw counts and the primary normalized rate are pre-committed bootstrap-disjoint at the $95\%$ level. The §5.2 reading \emph{(a)} initial formulation --- ``with $n=5$ agents the per-round price vector contains substantially more information than at $n=2$'' --- is empirically reversed: per-peer attention is denser at $n = 2$, not at $n = 5$. A refined reading \emph{(a$'$)} consistent with this anchor --- that the convergence-rate asymmetry at $n = 2$ reflects deeper per-peer attention density --- was operationally discriminated along three pre-committed channels: \emph{(i)} cross-capability scaling of the per-other-agent rate at $n = 2$ (Sonnet/Opus elevation supports \emph{(a$'$)}; regression toward GATE-5 supports the Haiku-class-artifact null) --- Stage 4; \emph{(ii)} the $Q$-learner positive control of \S\ref{sec:discussion-falsifiability}, whose absence of a chain-of-thought trace makes the basin-width gap presence/absence a no-CoT control on the attention-driven mechanism --- Stage 4; and \emph{(iii)} per-agent-round per-other-agent rate residualized on total CoT word count --- \emph{executed on Stage 2b data} via Addendum \#6 (paragraph below). Channels \emph{(i)} and \emph{(ii)} discriminate Haiku-class artifact vs cross-capability mechanism; channel \emph{(iii)} discriminates signal-density vs talk-volume scaling. None of the three independently adjudicates strategic-complexity content vs content-free reference noise --- a CoT-content-coding analysis is pre-committed for Stage 4 supplementary.

\paragraph{Talk-volume control --- Addendum \#6 (RFC 3161 stamped 2026-05-04).}
The §5.2.1 per-other-agent rate gap admits a candidate alternative explanation: GATE-2 might simply produce more total CoT text per round, so the per-other-agent reference count would scale mechanically with talk volume regardless of any per-peer attention-density mechanism. We pre-committed Addendum \#6 (RFC 3161 stamped 2026-05-04; SHA-256 \texttt{15d5663625a08660c27a9761ac2d722c1e124a75cc343aea728896a827b65606}) before any analyzer run, locking a pooled OLS regression of per-other-agent rate on total CoT word count, residualizing per record, and reporting per-condition residual means with bootstrap CIs at $B = 10{,}000$, fixed seed = 42. The talk-volume explanation is empirically ruled out: mean total word counts are nearly identical across conditions (GATE-5: $195.15$ words/CoT-round, $95\%$ CI $[194.44, 195.87]$; GATE-2: $194.67$, $[193.18, 196.16]$), so the per-other-agent rate gap cannot be attributed to one condition producing more CoT text. Residualizing the per-other-agent rate against the pooled OLS fit returns mean residuals of $-0.408$ ($[-0.416, -0.400]$) at GATE-5 and $+2.040$ ($[+1.951, +2.130]$) at GATE-2 --- disjoint $95\%$ bootstrap CIs with GATE-2 $\sim 5\times$ higher residual. The decision rule (Addendum \#6 \S B) fires \textsc{reversal survives residualization}; refined reading \emph{(a$'$)} is empirically supported on Stage 2b data even controlling for talk volume. Per-condition rates and residuals are reported in Table~\ref{tab:cot-residualized}.

\begin{table}[h]
\centering
\caption{Talk-volume control per Addendum \#6. Mean CoT word counts are nearly identical across conditions; the per-other-agent rate gap survives pooled OLS residualization with disjoint $95\%$ bootstrap CIs (GATE-2 mean residual $\sim 5\times$ higher than GATE-5). Bootstrap CIs reflect resampling variance only; see measurement validity caveat below.}
\label{tab:cot-residualized}
\begin{tabular}{lcc}
\toprule
Channel & GATE-5 ($n_{\text{rec}} = 7500$) & GATE-2 ($n_{\text{rec}} = 1500$) \\
\midrule
Mean CoT word count & $195.15$ $[194.44, 195.87]$ & $194.67$ $[193.18, 196.16]$ \\
Mean per-other-agent rate (Addendum \#5) & $0.508$ $[0.500, 0.516]$ & $2.954$ $[2.865, 3.045]$ \\
\textbf{Mean residual (talk-volume-controlled)} & $\mathbf{-0.408}$ $\mathbf{[-0.416, -0.400]}$ & $\mathbf{+2.040}$ $\mathbf{[+1.951, +2.130]}$ \\
Rate per $1000$ words (sensitivity) & $2.60$ $[2.56, 2.65]$ & $15.12$ $[14.66, 15.58]$ \\
\bottomrule
\end{tabular}
\end{table}

\paragraph{Statistical robustness sensitivity.}
The pooled OLS slope $b \approx 0.0045$ is dominated by GATE-5 records ($n = 7500$ vs $1500$); the within-GATE-2 OLS slope is $b_{n=2} \approx 0.0174$ ($\sim 3.9\times$ pooled), confirming that a condition-fixed-effects fit, FWL-equivalent to within-condition residualization, would shrink the GATE-2 mean residual rather than inflate it. The pooled spec is therefore directionally conservative for the headline claim. Round-level resampling does not cluster by repetition; under wild-cluster-by-rep, residual CI half-widths would inflate by approximately $2$--$3\times$, yielding post-inflation intervals (GATE-2 $\approx [+1.77, +2.31]$, GATE-5 $\approx [-0.42, -0.40]$) that remain disjoint. The $R^2 = 0.013$ establishes weak linear talk-volume sensitivity, not independence; nonlinearity at extreme word-counts is a residual threat bounded by the near-identical mean word counts ($\sim 195$) across conditions. The SR-M-1\ldots7 registry is sequentially-registered (Bonferroni-7 $\approx 0.0071$ if FWER control is desired), not formally FWER-corrected.

\paragraph{Q-learner control: necessity vs operativity (per Phys panel).}
The $Q$-learner positive control of \S\ref{sec:discussion-falsifiability} falsifies attention as \emph{necessary} for the basin-width gap; it does not falsify attention as \emph{operative} for LLMs. Both could be true simultaneously: a learning-rule-agnostic structural basin and an LLM-specific attention-density mechanism could coexist. Discriminating ``Haiku-class artifact $+$ structural basin'' from ``cross-capability-attention'' therefore requires the conjunction of channel \emph{(i)} (cross-capability scaling at $n = 2$) and channel \emph{(ii)} (Q-learner positive control), not either alone.

\paragraph{Measurement validity caveat.}
The cross-agent-reference rates reported here are produced by a regex-based analyzer over chain-of-thought text; bootstrap CIs reflect resampling variance only and do not bound construct validity. The token specifications in Addendum \#5 \S A.1--\S A.2 were pre-committed before any analyzer run, with one v1 $\to$ v2 widening (Addendum \#5 \S A.0) driven by inspection of one sample record per condition --- an asymmetric overfit to GATE-2 patterns is the specific threat that the deferred gold-standard subset would bound. We have published both analyzer sources (\texttt{pilot/analyze-cot-cross-references.ts}, \texttt{pilot/analyze-cot-residualized.ts}) and the per-condition output JSON files (\texttt{pilot/results/cot-cross-references-2026-05-03.json}, \texttt{pilot/results/cot-residualized-2026-05-04.json}), and we report both raw channels (literal, co-referent) alongside the per-other-agent normalized rate and the talk-volume-residualized rate. We have not produced a human-coded gold-standard subset with inter-annotator agreement, precision, and recall against the regex; the SciPy verifier (\texttt{pilot/verify-cot-cross-references.py}) is pre-committed but deferred. We treat the channel-(i) and channel-(ii) Stage 4 cross-model questions as exploratory open questions; the channel-(iii) talk-volume-controlled residualization is closed on Stage 2b data and reported above. The §5.2.1 conclusion is a directional reversal of reading \emph{(a)} as initially formulated, with refined reading \emph{(a$'$)} supported on Stage 2b data even under talk-volume control, pending Stage 4 cross-model generalization. Among the Stage 4 channels, the $Q$-learner positive control of \S\ref{sec:discussion-falsifiability} is a joint arbiter of refined reading \emph{(a$'$)} alongside channel \emph{(i)} cross-capability scaling: its absence of any chain-of-thought trace makes the basin-width gap presence/absence under matched parameters a no-CoT control that adjudicates whether attention is \emph{necessary} for the basin gap (consistent with the necessity-vs-operativity decomposition above).

\subsection{Implications for the field}
\label{sec:discussion-implications}

Stage 2b reports three findings about LLM-agent Bertrand pricing --- at distinct
evidentiary tiers --- each grounded in the locked $n_{\text{rep}} = 15$ +
$n_{\text{rep}} = 30$ dataset and replicated by an independent SciPy verifier.
\emph{Formally established at GATE-5} ($n_{\text{rep}} = 30$): LLM agents \emph{at
$n = 5$} converge to a bimodal terminal-state distribution ($k = 2$ recovered in $7/10$
gap-statistic resamples; $30/30$ basin-stable; $\Delta_{\text{profit}} = 0.49$ at
convergence, capturing roughly half the Nash-to-Monopoly profit gap).
\emph{Descriptively present at GATE-2} ($n_{\text{rep}} = 15$): the same
$\Delta_{\text{profit}}$ landing point ($0.45$) is reached through structurally
different convergence dynamics ($1/15$ strict-converged, $12/15$ basin-stable with
Wilson CI $[0.548, 0.930]$ that overlaps the $0.80$ PROCEED threshold) --- the
asymmetry that motivates the Stage 3 design but does not on its own discriminate
basin-width-by-agent-count from horizon-and-sample-size effects
(\S\ref{sec:discussion-falsifiability}). \emph{Conditional, multi-condition
characterization}: extending \citet{Cao_2026_llm_collusion}'s LLM bistability finding,
the basin signature is recovered at $n = 5$ and not recovered at $n = 2$ within Stage
2b power (\S\ref{sec:limitations-envelope}). These outcomes are partial-coordination
--- supra-competitive without explicit agreement --- and \emph{track} qualitatively
the regulator-relevant case identified by \citet{ezrachi2016virtual,asker2021ai_pricing}.
We do not quantify the consumer-surplus loss implied by $\Delta_{\text{profit}} \approx 0.45$--$0.49$
at the Calvano logit-demand calibration here; the welfare-loss translation is deferred to
companion IO work building on Stage 3.

\paragraph{Antitrust and regulatory framing.}
The basin-width-by-agent-count reading aligns the methodological frame with the regulator-relevant case. \citet{ezrachi2016virtual}'s Predictable Agent category is most acute at small agent counts --- duopolies and tight oligopolies, where observability of explicit coordination is hardest, are precisely where the wider-basin reading predicts the longest, noisiest convergence dynamics. A regulator inheriting our framework looks at convergence dynamics, not just terminal prices: the population of long-transient duopolies that have not statically converged but are basin-stable in the wider sense is the regulator-relevant case, and the basin lens makes that population legible. The methodological contribution is that we report it as a separate, pre-registered secondary verdict rather than absorbing it into the primary convergence statistic.

\paragraph{Multi-agent LLM safety.}
The basin-of-attraction lens is method-importable to other multi-agent LLM settings beyond pricing --- negotiation, contracting, market-making, allocation. \citet{Hammond_2025_multiagent_risks,Riedl_2025_emergent} survey emergent multi-agent LLM behaviors but generally treat outcomes as point estimates rather than as basin samples. Reporting cross-rep distributions plus basin-stability secondary verdicts plus a $k$-selection diagnostic is a more honest accounting of multi-agent LLM behavior than reporting a single expected outcome, and we recommend the practice as a default for multi-agent LLM empirical work.

\paragraph{LLM-versus-$Q$-learner family generalization.}
Our results suggest the Calvano-Bertrand basin signature is robust across algorithm families. $Q$-learners with explicit reward signals \citep{calvano_etal_2020_ai_algorithmic_pricing_collusion}, Cournot $Q$-learners \citep{waltman_kaymak_2008_qlearning_cournot}, and now LLM agents with no reward signal at all all converge to multimodal terminal distributions in repeated pricing games. This is consistent with \citet{asker2021ai_pricing}'s framing that algorithmic collusion is a property of the learning rule and game structure, not of any specific algorithmic family. Stage 4 cross-model replication (Haiku sizes, Sonnet, Opus, non-Anthropic frontier models) is designed to test this generalization formally.

\subsection{Falsifiability and what would change our conclusions}
\label{sec:discussion-falsifiability}

We pre-commit four specific falsification conditions for Stage 3 and Stage 4. First, the basin-width-by-agent-count hypothesis is rejected if Stage 3 ($n_{\text{rep}} \geq 30$, horizon $\geq 100$) produces statistically indistinguishable strict-convergence rates at $n = 2$ and $n = 5$ (Wilson 95\% CIs on the per-condition convergence rates overlapping by more than half their width). Second, the asymmetry headline is rejected if Stage 3 $\Delta_{\text{profit}}$ at $n = 2$ falls outside $[0.30, 0.60]$, the interval bracketing the Stage 2b point estimate $\pm$ the bootstrap CI half-width. Third, the LLM-versus-$Q$-learner family-generalization claim is rejected if a parallel $Q$-learner experiment under the same Bertrand setting at $n \in \{2, 5\}$ does not exhibit a similar basin-width asymmetry; this $Q$-learner positive-control experiment is pre-committed in \texttt{pilot/VALIDATION\_PROTOCOL.md} (15 runs of an exact Calvano $n=2$ replication) and runs alongside the Stage 4 cross-LLM-model replication. Fourth, the basin-width-by-agent-count hypothesis predicts a monotonic relationship between agent count and the strict-convergence rate; Stage 3 will report the per-agent-count strict-convergence rate at $n \in \{2, 3, 4, 5\}$, and the hypothesis is rejected if the curve is non-monotonic or if the $n = 3$ and $n = 4$ rates are not bracketed by the $n = 2$ and $n = 5$ rates within Wilson 95\% CIs. Each condition is locked here in advance of Stage 3 data collection.

\subsection{Wu et al.\ 2025 simplicity-bias mechanism (locked verbatim per Addendum \#1 \S D)}
\label{sec:discussion-wu2025}

We pre-committed (Addendum \#1 \S D, RFC 3161 stamped 2026-04-23) to the \citet{wu2025whenMoreIsLess} simplicity-bias mechanism as the candidate explanation for any negative correlation between chain-of-thought length and basin welfare loss: if the cooperative-pricing answer is representationally simpler for a Haiku-class pretrained agent, simplicity-bias predicts shorter reasoning at more-cooperative basins. The point-estimate signs at both agent counts are consistent with this prediction ($r(\text{reasoning length}, \Delta_{\text{profit}}) = -0.30$ at $n = 5$ and $-0.47$ at $n = 2$ over the locked rounds 41--50 final-window aggregation), but Fisher-$z$ 95\% confidence intervals span zero in both conditions ($[-0.60, +0.06]$ at $n = 5$; $[-0.79, +0.05]$ at $n = 2$); $n_{\text{rep}} = 15$ at GATE-2 caps discrimination here, see \S\ref{sec:limitations-envelope} for the cross-result envelope. The pre-registered falsifiability clause --- $r(\text{reasoning length}, \Delta_{\text{profit}}) > 0$ at $n = 5$ falsifies simplicity-bias as the operative mechanism --- is NOT triggered: the $n = 5$ point estimate is negative. The mechanism is consistent with the data direction but not strongly confirmed at this sample size. The pre-committed Stage 4 cross-model replication (across Haiku sizes, Sonnet, Opus, and non-Anthropic frontier models) is the formal falsification test.
